\documentclass[11pt]{article}

% Use the official NeurIPS style if it is available locally. Fall back to a
% simple article layout so the manuscript source remains buildable in a bare
% environment.
\newif\ifneuripsstyle
\IfFileExists{neurips_2024.sty}{
  \neuripsstyletrue
  \usepackage[final]{neurips_2024}
}{
  \neuripsstylefalse
  \usepackage[margin=1in]{geometry}
  \usepackage[numbers]{natbib}
  \usepackage[colorlinks=true,allcolors=blue]{hyperref}
}

\usepackage[utf8]{inputenc}
\usepackage[T1]{fontenc}
\usepackage{microtype}
\usepackage{amsmath}
\usepackage{amssymb}
\usepackage{booktabs}
\usepackage{array}
\usepackage{graphicx}
\usepackage{multirow}
\usepackage{url}

\title{Minimal Execution Interventions for Post-Hoc Debugging of Tool-Using Agents}

\author{Anonymous Submission}

\begin{document}

\maketitle

\begin{abstract}
Tool-using language agents often fail because a single intermediate action
corrupts the downstream execution state, yet standard benchmark evaluation
typically exposes only the final reward and a raw trajectory trace. We propose
post-hoc minimal patch recovery, a debugging methodology that treats failed
agent trajectories as replayable programs. Given a failed run and full
checkpointed environment state, our method restores an intermediate snapshot,
applies a small structured intervention, replays execution, and tests whether
the benchmark outcome flips from failure to success. The smallest successful
intervention serves two purposes at once: it is a repair and a causal autopsy
of the root failure. We instantiate this framework on real tau-style agent
environments with deterministic state, tool side effects, and benchmark
evaluators. The current system supports multiple patch families, including tool
argument edits, tool-call replacement, insertion, deletion, context edits, and
bounded continuation replacement, together with explicit search-cost
accounting. On controlled synthetic corruptions, the prototype recovers
failures end to end and reports exact fault-localization metrics such as top-k
accuracy and mean reciprocal rank. On imported natural failures from saved
benchmark runs, the current prototype already demonstrates strict non-oracle
recovery without oracle suffix reuse in both single-case and bundle settings.
A dedicated retail identity-lookup failure is repaired via staged continuation
from replayed state, and a clean multi-model mixed-domain bundle of $32$
replay-failing natural cases yields $13/32$ strict recoveries under bounded
deterministic search, with per-domain recovery $\textit{airline}=10/16$ and
$\textit{retail}=3/16$. These results position replayable execution
intervention as a promising methodology for agent debugging: more causal than
prompt-only retry, more interpretable than unconstrained regeneration, and
naturally compatible with benchmark-native evaluation.
\end{abstract}

\section{Introduction}

Language agents that interact with tools, databases, websites, or code
repositories fail for reasons that are often difficult to diagnose from final
task reward alone. In a typical benchmark run, we observe a user goal, a
multi-step trajectory, and a binary or scalar outcome. If the agent fails, the
default response is usually to inspect the raw trace, retry with a stronger
model, or ask the model to reflect and regenerate. These strategies can improve
performance, but they are weak debugging tools: they do not cleanly isolate
which intermediate decision actually caused failure, and they do not
distinguish whether the real bottleneck is action selection, stale context,
harmful state mutation, or continuation from a corrupted state.

This paper argues that replayable agent environments enable a stronger form of
post-hoc analysis. If the environment state can be checkpointed and restored,
then a failed trajectory can be treated as a stateful program rather than a
frozen transcript. We can pause execution at an intermediate step, edit a small
piece of state or a single action, replay from that fork, and test whether the
benchmark reward changes. This makes it possible to ask a causal question:
which minimal intervention is sufficient to flip the final outcome?

We study this question through post-hoc minimal patch recovery. The goal is not
to train a new agent or to win a benchmark with a larger model. Instead, the
goal is to turn failed trajectories into analyzable objects. A successful patch
provides both a repair and an explanation: it identifies a concrete failure
site, measures the cost of finding a fix, and yields an autopsy artifact that
can be aggregated across many failures.

Our current contribution is methodological rather than leaderboard-oriented. We
introduce a benchmark-native debugging protocol that combines checkpointed
replay, structured intervention families, bounded search, and reward-based
evaluation. Its main outputs are success recovery rate, patch minimality,
search cost, localization quality, and structured autopsy reports.

\section{Problem Setup}

We model a benchmark task as an interactive environment $E$ with an initial
task specification $x$, an environment state $s_0$, a transition function over
tool calls and user messages, and an evaluator $R$ that returns the benchmark
outcome for a completed trajectory. Executing a baseline policy $\pi$ produces
a trajectory
\[
\tau = (s_0, a_0, o_0, s_1, a_1, o_1, \ldots, s_T, y),
\]
where $a_t$ is the agent action at step $t$, $o_t$ is the resulting tool or
environment observation, $s_t$ is the pre-step environment state, and $y$ is
the final task outcome.

We assume a failed trajectory with $R(\tau)=0$. A patch candidate is a tuple
\[
p = (t, f, \delta),
\]
where $t$ is the target step, $f$ is a patch family, and $\delta$ is the patch
payload. Applying $p$ means:
\begin{enumerate}
  \item restore the saved pre-step snapshot $s_t$;
  \item modify the action, context, or execution state according to $\delta$;
  \item resume execution under a specified continuation policy;
  \item evaluate the patched rollout with the original benchmark evaluator.
\end{enumerate}

Let $\tau'(p)$ be the replayed rollout after applying patch $p$. The search
problem is:
\[
\text{find } p^\star \text{ such that } R(\tau'(p^\star)) = 1
\]
while minimizing patch size subject to a bounded search budget. In our
implementation, the budget can include localization cost, patch-proposal cost,
rollout count, and model-token usage.

Two continuation settings are important:
\begin{itemize}
  \item \textbf{Strict replay}: continue from the replayed state without using
  oracle future actions from the benchmark trace.
  \item \textbf{Oracle continuation}: use the benchmark reference suffix only
  as an upper-bound analysis tool.
\end{itemize}

This distinction matters because it separates two different bottlenecks:
localizing the faulty step and generating a good continuation after that step
is changed.

\section{Method}

\subsection{Failure Collection}

We construct two classes of corpora:
\begin{itemize}
  \item \textbf{Synthetic corruption failures}, created by perturbing
  otherwise-valid benchmark-shaped trajectories so that the true fault step is
  known.
  \item \textbf{Natural imported failures}, loaded from saved tau-style result
  files produced by an external baseline agent.
\end{itemize}

For every trajectory we store prompt context, action, tool result, pre-step
snapshot, post-step snapshot, and final outcome metadata.

\subsection{Fault Localization}

Given a failed trajectory, the search ranks candidate intervention steps. The
current system supports deterministic ranking strategies such as heuristic,
reverse, chronological, latest-only, oracle-fault-step for synthetic analysis,
and random candidate order. The point of localization is not only speed. It
also enables direct measurement of root-cause accuracy on synthetic failures
with known fault steps.

\subsection{Patch Proposal}

The current intervention families are:
\begin{itemize}
  \item \texttt{tool\_args}
  \item \texttt{tool\_call\_replace}
  \item \texttt{tool\_insertion}
  \item \texttt{tool\_deletion}
  \item \texttt{context\_edit}
  \item \texttt{continuation\_replace}
  \item \texttt{tool\_call\_with\_oracle\_suffix}
\end{itemize}

These patch families operate on structured execution objects rather than
unconstrained free-form text. That keeps search interpretable and makes patch
size measurable by family-specific criteria.

\subsection{Replay and Continuation}

For each candidate patch we restore the saved pre-step snapshot, apply the
patch, execute the modified step when appropriate, and resume execution. The
continuation policy is a critical part of the method. Some failures can be
recovered by editing one action and replaying the original suffix. Others
cannot, because the patched state changes later observations and invalidates
the old suffix. The current system therefore distinguishes simple strict
replay, staged continuation from the replayed state, and oracle suffix
continuation as an upper bound only.

\subsection{Selection and Autopsy Generation}

Among successful patches, the system selects the smallest one, breaking ties by
lower search cost. It then emits a structured autopsy containing the failure
identifier and task metadata, the root-cause step, patch family, patch size,
original versus patched action or state fragment, continuation mode, total
search cost, and a short natural-language explanation.

\section{Evaluation Protocol}

The evaluation story has two tracks that remain separate in the manuscript.

\subsection{Synthetic Evaluation}

Synthetic corruptions provide known fault steps and known pre-corruption
reference behavior. This enables exact localization accuracy, minimality
analysis under known single-fault settings, and strategy comparisons where
recovery and localization can diverge.

\subsection{Natural Failure Evaluation}

Natural failures are imported from saved benchmark result files. These runs are
ecologically more realistic because the failures were produced by an external
baseline agent rather than injected by our pipeline. The main natural headline
metric is strict non-oracle recovery. Oracle continuation is reported
separately as an upper bound and is never mixed into the main claim.

\subsection{Core Metrics}

The paper tracks success recovery rate, patch size, total search cost,
success@budget, per-family recovery, per-domain recovery, localization metrics
on synthetic failures, and the strict-versus-oracle continuation gap.

\section{Preliminary Results}

Table~\ref{tab:synthetic} summarizes the current controlled synthetic evidence.
The main lesson is that recovery and localization should be measured
separately. A strategy can recover some failures without accurately isolating
the causal step.

\begin{table}[t]
  \centering
  \caption{Current controlled synthetic strategy results from
  \texttt{artifacts/synthetic\_strategy\_tables\_3/paper\_tables.json}.}
  \label{tab:synthetic}
  \begin{tabular}{lcccc}
    \toprule
    Strategy & Recovery & Top-1 Loc. & MRR & Token Cost \\
    \midrule
    Heuristic & $3/3$ & $1.00$ & $1.00$ & $3309$ \\
    Latest only & $3/3$ & $1.00$ & $1.00$ & $3165$ \\
    Chronological & $2/3$ & $0.00$ & $0.164$ & $2109$ \\
    Random candidate & $2/3$ & $0.00$ & $0.264$ & $1961$ \\
    No repair & $0/3$ & -- & -- & $0$ \\
    \bottomrule
  \end{tabular}
\end{table}

Table~\ref{tab:natural} summarizes the strongest current natural-failure
evidence. The project is no longer only a synthetic toy story: one retail case
is recovered by staged continuation after fixing a bad lookup step, and a clean
multi-model mixed-domain bundle reaches $13/32$ strict recoveries. At the same
time, the broader retail natural import still shows that naive local repair is
insufficient on many real failures.

\begin{table}[t]
  \centering
  \caption{Current natural-failure evidence from saved repository artifacts.}
  \label{tab:natural}
  \begin{tabular}{p{4.2cm}cccc}
    \toprule
    Setting & Recovery & Cost & Main Family & Artifact \\
    \midrule
    Broader retail strict replay & $0/3$ & $280$ & none &
    \texttt{natural\_search} \\
    Retail staged continuation & $1/1$ & $483$ &
    \texttt{continuation\_replace} &
    \texttt{natural\_runtime\_smoke\_search} \\
    Mixed-domain strict bundle & $13/32$ & $2000$ &
    \texttt{tool\_deletion} &
    \texttt{paper\_bundle\_multimodel\_32} \\
    Mixed-domain oracle upper bound & $13/32$ & $2000$ &
    \texttt{tool\_deletion} &
    \texttt{paper\_bundle\_multimodel\_32/oracle} \\
    \bottomrule
  \end{tabular}
\end{table}

The mixed-domain bundle contains failures imported from runs originally produced
by \texttt{claude-3-7-sonnet-20250219}, \texttt{gpt-4.1-2025-04-14},
\texttt{gpt-4.1-mini-2025-04-14}, and \texttt{o4-mini-2025-04-16}. Under the
current bounded deterministic search configuration, per-domain strict recovery
is $10/16$ for airline and $3/16$ for retail. These results are encouraging,
but they remain smoke-scale rather than benchmark-scale.

\section{Limitations and Scope}

The current project is promising but not yet at final submission scale. First,
the natural-failure evidence is real but still modest. The repository has
verified strict non-oracle recoveries on saved benchmark failures, including a
clean thirty-two-case mixed-domain bundle, but not yet at the scale needed for
a strong benchmark-wide claim.

Second, continuation remains a major bottleneck. Some failures are easy to
localize yet hard to recover because the patched action changes the future
state distribution. This is exactly why the strict-versus-oracle gap is
informative, but it also limits current natural recovery.

Third, several intervention families exist in the interface before they are
fully mature as high-quality proposers across domains. The paper should be
careful to distinguish interface support from fully competitive evaluated
methods.

Fourth, autopsy usefulness has not yet been validated with human studies or a
large expert annotation pass. At the current stage, autopsies should be treated
as structured artifacts and qualitative evidence, not yet as fully validated
human-facing explanations.

Finally, the paper is about debugging methodology, not deployment. The current
system operates on offline benchmark trajectories and replayable environments.
It is not a claim about safe online automated intervention in real customer
service systems.

\section{Related Work}

\subsection{Interactive Agent Benchmarks}

Recent work on agent evaluation has established the importance of realistic,
interactive benchmarks with tool use, persistent state, and goal-conditioned
tasks. Tau-style customer-service environments, WebArena, WorkArena,
AgentDojo, and software-oriented settings such as SWE-bench all share a core
insight: final task success depends on multi-step interaction with an external
environment, not just single-turn text quality
\citep{tau_style,web_arena,work_arena,agent_dojo,swe_bench}.

Our work is complementary to this benchmark line. These benchmarks tell us
whether an agent succeeded or failed, and they often provide rich traces, but
they do not by themselves offer a principled way to intervene on failed runs.
We build on benchmark-style environments rather than replacing them. The key
difference is methodological: instead of treating a failed trajectory as a dead
artifact, we treat it as a replayable object that can be restored, patched, and
re-evaluated under the original benchmark reward.

\subsection{Reflection, Retry, and Self-Repair}

A second relevant line of work studies how language models can critique and
improve their own outputs through reflection, refinement, or search. This
includes approaches in the spirit of Reflexion, Self-Refine, tree-search style
deliberation, and self-debugging pipelines
\citep{reflexion,self_refine,tree_of_thoughts,self_debugging}.

Our setting is related but importantly different. Reflection-style methods
typically operate in prompt space: they ask the model to reconsider, critique,
or regenerate a solution. In contrast, we operate in execution space. We do
not only ask the model to try again; we restore the exact environment state at
a chosen step, apply a targeted structured intervention, and measure whether
the real downstream environment reward changes. This yields a more causal notion
of repair than prompt-only retry, while also exposing a strict-versus-oracle
continuation gap that can be studied directly.

\subsection{Fault Localization, Program Repair, and Counterfactual Explanations}

Our problem is also closely related to classic ideas from fault localization,
delta debugging, automated program repair, and minimal counterfactual
explanations \citep{delta_debugging,program_repair,minimal_counterfactuals}.
In those settings, one seeks a small change that either explains or repairs
undesirable behavior. The smallest successful edit is often more informative
than a large unconstrained rewrite.

We adopt the same spirit in the agent setting, but the object being edited is
not source code alone. It is a trajectory embedded in a live task environment
with dialogue history, retrieved context, tool calls, and stateful side
effects. That makes the setting harder than ordinary static repair: changing an
action can alter later observations, invalidate the original suffix, or require
fresh continuation planning from the replayed state.

\section{Conclusion}

This project supports a simple but powerful thesis: failed agent trajectories
in replayable environments can be debugged through causal execution
intervention rather than only through prompt-level retry. By restoring a saved
state, applying a bounded structured patch, and re-evaluating the resulting
rollout in the original benchmark environment, we can convert opaque failures
into concrete autopsies with explicit recovery cost.

The current repository already validates the main shape of that thesis.
Controlled synthetic failures can be localized and repaired end to end.
Natural saved benchmark failures are harder, but they already show the most
important methodological pattern: some failures are irrecoverable under naive
single-step replay, some become recoverable under staged continuation, and some
can be recovered strictly once the intervention space includes richer repair
operations such as action deletion.

The most important next step is scale. A final paper should move from smoke
artifacts to a frozen corpus of hundreds of natural failures across at least
\textit{retail} and \textit{airline}, while preserving the current split
between synthetic localization analysis and natural ecological evaluation. The
strongest version of the paper will also strengthen strict continuation, add
model-backed proposal ablations, replace placeholder citations with exact
bibliographic entries, and produce camera-ready figures directly from saved
JSON artifacts.

\bibliographystyle{plainnat}
\bibliography{references}

\end{document}
